\subsection{POINCAR\'E-BIRKHOFF-WITT 定理的整形式}

令 g 为有限维 Q-李代数，\(U(\mathfrak{g})\) 为其包络双代数。若 I 为全序集，\(\mathbf{x} = (x_i)_{i \in \mathrm{I}}\) 为 g 中元素的一个族，\(\mathbf{n} = (n_i)_{i \in \mathrm{I}} \in \mathbf{N}^{(\mathrm{I})}\) 为一个多重指标，置

\[
\mathbf {x} ^ {(\mathbf {n})} = \prod_ {i \in \mathrm{I}} \frac {x _ {i} ^ {n _ {i}}}{n _ {i} !},\tag{5}
\]

其中乘积在 \(\mathrm{U}(\mathfrak{g})\) 中按照全序集 I 计算。

定理 1。令 \(\mathcal{U}\) 为双代数 \(\mathrm{U}(\mathfrak{g})\) 中的双序。令 \(\mathcal{G} = \mathcal{U} \cap \mathfrak{g}\)，这是李代数 \(\mathfrak{g}\) 中的一个序。令 \((x_i)_{i \in \mathrm{I}}\) 为 \(\mathcal{G}\) 的一个基。给 I 赋予一个全序，并假设对每个 \(\mathbf{n} \in \mathbf{N}^{\mathrm{I}}\)，给定 { 中的一个元素 }[{n}]\(\mathcal{U}\)，使 {[}n{]} - \(\mathbf{x}^{(\mathbf{n})}\) 在 \(< |\mathbf{n}|\) 中的滤子次数小于 \(\mathrm{U}(\mathfrak{g})\)。则 { 的 }[{n}]\(\mathbf{n} \in \mathbf{N}^{\mathrm{I}}\) 所成的族是 Z-模 \(\mathcal{U}\) 的一个基。

对 \(p \in N\)，令 \(U_{p}(\mathfrak{g})\) 为 \(U(\mathfrak{g})\) 中滤子次数 \(\leq p\) 的元素所成的集合；则 \(U_{p}(\mathfrak{g}) / U_{p-1}(\mathfrak{g})\) 的 \(\mathbf{x}^{(\mathbf{n})}\) 在 \(|n| = p\) 中的像构成此 Q-向量空间的一个基（第 I 章第2节第7号定理1）；因此 {[}n{]} 构成 Q-向量空间 \(U(\mathfrak{g})\) 的一个基。还需证明以下断言（其中置 \(M = N^{I}\)）：

\((^{*})\) 若 \(u\in \mathcal{U}\)、\((a_{\mathbf{n}})\in \mathbf{Z}^{(\mathrm{M})}\) 且 \(d\in \mathbf{N} - \{0\}\) 满足

\[
d u = \sum_ {\mathbf {n} \in \mathrm{M}} a _ {\mathbf {n}} [ \mathbf {n} ],\tag{6}
\]

则 \(d\) 整除每个 \(a_{\mathbf{n}}\)。

对每个整数 \(r \geq 0\)，引入迭代余积

\[
c _ {i}: \mathcal {U} \rightarrow \mathbf {T} ^ {r} (\mathcal {U}) = \mathcal {U} \otimes \mathcal {U} \otimes \dots \otimes \mathcal {U};
\]

按定义，\(c_0\) 是 \(\mathcal{U}\) 的余单位，\(c_1 = \mathrm{Id}_{\mathcal{U}}\)，\(c_2 = c\)（\(\mathcal{U}\) 的余积），并且当 \(r \geq 2\) 时，\(c_{r+1}\) 定义为复合映射 \(p \circ (c_r \otimes 1) \circ c\)：

\[
\mathcal {U} \stackrel {{c}} {{\longrightarrow}} \mathcal {U} \otimes_ {\mathbf {Z}} \mathcal {U} \stackrel {{c _ {r} \otimes 1}} {{\longrightarrow}} \mathbf {T} ^ {r} (\mathcal {U}) \otimes_ {\mathbf {Z}} \mathcal {U} \stackrel {{p}} {{\longrightarrow}} \mathbf {T} ^ {r + 1} (\mathcal {U})
\]

其中 p 由利用代数 \(\mathbf{T}(\mathcal{U})\) 中的乘法来定义。此外，考虑 U 到 \(\pi\) 的典则投影 \(U^{+} = Ker c_{0}\)，以及复合映射

\[
c _ {r} ^ {+} = \mathbf {T} ^ {r} (\pi) \circ c _ {r}: \mathcal {U} \to \mathbf {T} ^ {r} (\mathcal {U} ^ {+}).
\]

引理 1。令 \(\mathbf{n} \in \mathbf{N}^{\mathrm{I}}\)。若 \(|\mathbf{n}| < r\)，则 \(c_{r}^{+}([ \mathbf{n}]) = 0\)。若 \(|\mathbf{n}| = r\)，则

\[
c _ {r} ^ {+} ([ \mathbf {n} ]) = \sum_ {\varphi} x _ {\varphi (1)} \otimes x _ {\varphi (2)} \otimes \dots \otimes x _ {\varphi (r)},\tag{7}
\]

其中 \(\varphi\) 属于从 \(\{1,2,\ldots,r\}\) 到 I 的映射所成的集合，并且对每个值 \(i\in I\)，它取值 i 恰好 \(n_{i}\) 次。

由命题1，

\[
c _ {r} \left(\mathbf {x} ^ {(\mathbf {n})}\right) = \sum \mathbf {x} ^ {\left(\mathbf {p} _ {1}\right)} \otimes \dots \otimes \mathbf {x} ^ {\left(\mathbf {p} _ {r}\right)}
\]

其中求和遍历由 M 中 r 个元素组成的序列 \((\mathbf{p}_{1},\ldots,\mathbf{p}_{r})\) 所成的集合，这些元素满足 \(p_{1}+\cdots+p_{r}=n\)。根据第 II 章第1节第3号命题6，将 \(c_{r}^{+}\) 通过线性延拓为从 \(\mathrm{U}(\mathfrak{g})\) 到 \(\mathbf{T}^{r}(\mathrm{U}^{+}(\mathfrak{g}))\) 的映射后，它在 \(\mathrm{U}_{r-1}(\mathfrak{g})\) 上为零。因此，对 \(r\geq|\mathbf{n}|\)，

\[
c _ {r} ^ {+} ([ \mathbf {n} ]) = c _ {r} ^ {+} (\mathbf {x} ^ {(\mathbf {n})}) = \sum \pi (\mathbf {x} ^ {(\mathbf {p} _ {1})}) \otimes \dots \otimes \pi (\mathbf {x} ^ {(\mathbf {p} _ {r})}).\tag{8}
\]

当 \(r > |n|\) 时，关系 \(p_{1} + \cdots + p_{r} = n\) 蕴含至少有一个 \(p_{i}\) 为零，所以 \(c_{r}^{+}([n]) = 0\)。当 \(r = |n|\) 时，(8) 的第三项中非零的项只有满足 \(|p_{1}| = \cdots = |p_{r}| = 1\) 的项，因而得到 (7)。

回到定理1的证明。沿用 (*) 的记号，并对 \(|n|\) 作降阶归纳，证明 d 整除 \(a_{n}\)，当 \(|n|\) 充分大时这显然成立。若对 \(a_{n}\) 有 d 整除 \(|n| > r\)，则置

\[
u ^ {\prime} = u - \sum_ {| \mathbf {n} | > r} (a _ {\mathbf {n}} / d) [ \mathbf {n} ] \in \mathcal {U},
\]

有

\[
d u ^ {\prime} = \sum_ {| \mathbf {n} | \leq r} a _ {\mathbf {n}} [ \mathbf {n} ].\tag{9}
\]

对从 \(\varphi\) 到 I 的任意映射 \(\{1,2,\ldots,r\}\)，置

\[
e _ {\varphi} = x _ {\varphi (1)} \otimes \dots \otimes x _ {\varphi (r)}
\]

并令 \(a_{\varphi} = a_{\mathbf{n}}\)，其中 \(\mathbf{n} = (\operatorname{Card}\varphi^{-1}(i))_{i\in \mathrm{I}}\)。由引理1，(9) 蕴含

\[
d c _ {r} ^ {+} (u ^ {\prime}) = \sum_ {\varphi \in \mathrm{I} ^ {r}} a _ {\varphi} e _ {\varphi}\tag{10}
\]

所以 \(c_r^+(u') \in \mathbf{T}^r(\mathcal{U}^+) \cap \mathbf{Q}\mathbf{T}^r(\mathcal{G})\)。但是 \(\mathcal{G}\) 作为 \(\mathcal{U}^+\) 的子模是一个直因子（《代数》第 VII 章第4节第3号定理1的推论），所以子模 \(\mathbf{T}^r(\mathcal{G})\) 是 \(\mathbf{T}^r(\mathcal{U}^+)\) 的直因子，从而 \(c_r^+(u') \in \mathbf{T}^r(\mathcal{G})\)。另一方面，由假设，\(x_i\) 构成 \(\mathcal{G}\) 的一个基，所以 \(e_\varphi\) 构成 \(\mathbf{T}^r(\mathcal{G})\) 的一个基。于是 (10) 证明 d 整除 \(d\)\(a_\varphi\)，即当  时 d 整除 \(a_{\mathbf{n}}\)。这证明了 (*)。\(|\mathbf{n}| = r\)
% semantic-fragments: legacy-input-seam
\relax{}
